\section{结语}
\label{sec:conclusion}

$C_n(1,2)$ 上的符号谱问题由圈 flux 而非单条边符号自然控制。负半胞通量给出一般单项式
chiral 对合，并把平方 Bloch 问题维数减半。周期八是该机制首次产生平方谱边小于 $8$ 的
本原周期，相应的对径双缺陷轨道在该周期内唯一。它的额外 block 结构使全部色散可以显式
求出，因此两种有限 holonomy 都有精确谱半径；正 holonomy 随后给出严格优于 twisted
benchmark 的无限图族。

精确值 $\sqrt\etaedge$ 是本文显式符号达到的谱半径，但尚未证明它等于
$m(C_{8L}(1,2))$。剩余问题可以精确表述为：是否存在一个无限 $L$ 子族满足
\[
 m(C_{8L}(1,2))=\sqrt\etaedge,
\]
以及在等号成立时如何刻画全部最小 switching 类。半胞判据和缺陷约束为这一问题提供
两组结构坐标；本文不需要对所有偶数阶或全部周期相作额外声明。
